% Rozdzial 20. Fragment do \input; bez preambuly.
% Wymaga listings, enumitem, tabularx, booktabs, hyperref i makr
% \code, \term, \nota. Bazowo C++20.
\chapter{Projektowanie bibliotek i nowoczesnych interfejsów API}
\label{chap:cpp20-api}

\section{Czym jest dobry interfejs biblioteki?}

Biblioteka to nie tylko poprawny algorytm, lecz przede wszystkim
\emph{umowa} między jej autorem a użytkownikiem. Umowa obejmuje typy,
sygnatury, ograniczenia, własność zasobów, czas życia, obsługę błędów,
bezpieczeństwo wyjątkowe, własności współbieżne i kompatybilność wersji.
Każda decyzja ujawniona w publicznym nagłówku może stać się
kosztownym zobowiązaniem na lata. Celem jest interfejs, w którym
\emph{poprawne użycie jest naturalne, a błędne możliwie trudne}.

\begin{center}\small
\begin{tabularx}{\textwidth}{@{}p{0.30\textwidth}X@{}}
\toprule
Własność API & Pytanie kontrolne \\
\midrule
Kontrakt & Co jest dozwolonym wejściem, a co wynikiem? \\
Własność & Kto posiada obiekt, kto go zwalnia i kiedy? \\
Bezpieczeństwo & Czy błędne wejście może spowodować UB albo wyciek? \\
Wyjątki & Jakie gwarancje pozostają po niepowodzeniu? \\
Współbieżność & Co wolno wywoływać równocześnie i co jest synchronizowane? \\
Wydajność & Jakie kopiowania, alokacje i synchronizacje są narzucane? \\
Kompatybilność & Które zmiany są zgodne źródłowo i binarnie? \\
\bottomrule
\end{tabularx}
\end{center}

\nota{Rozdziały o smart pointerach, concepts, type erasure, wyjątkach,
linkowaniu i ABI dostarczają narzędzi. Tu rozpatrujemy je razem, jako
wybory projektowe podejmowane przy budowaniu stabilnych bibliotek.}

\section{Granica interfejsu: dane, widoki i własność}
\label{sec:cpp20-ownership}

Najpierw wybieramy, czy argument jest \emph{wejściem}, \emph{wyjściem},
\emph{stanem współdzielonym} czy \emph{zasobem przekazywanym na własność}.
Dopiero potem wybieramy konkretny typ C++.

\begin{center}\small
\begin{tabularx}{\textwidth}{@{}p{0.37\textwidth}X@{}}
\toprule
Parametr & Interpretacja w typowym interfejsie \\
\midrule
\code{T} & Wartość własna; przekazanie lub przejęcie kopii/przeniesienia. \\
\texttt{const T\&} & Pożyczony, niezmieniany obiekt, który musi przeżyć wywołanie. \\
\texttt{T\&} & Obowiązkowy, modyfikowalny obiekt pożyczony. \\
\code{T*} & Zwykle opcjonalny, niewłasnościowy dostęp; dopuszcza \code{nullptr}. \\
\code{unique_ptr<T>} & Przeniesienie lub przechowywanie wyłącznej własności. \\
\code{shared_ptr<T>} & Współdzielona własność, nie sam odczyt obiektu. \\
\code{span<T>} & Pożyczony widok ciągłego bufora, bez informacji o własności. \\
\code{string_view} & Pożyczony widok tekstu, możliwy problem czasu życia. \\
\bottomrule
\end{tabularx}
\end{center}

\subsection{Parametry i brak zbędnej alokacji}

\begin{lstlisting}[caption={Interfejs operujacy na dowolnym ciaglym buforze (C++20).}]
#include <cstddef>
#include <span>
#include <stdexcept>

// Nie przechowuje span po powrocie z funkcji.
double weighted_sum(std::span<const double> a,
                    std::span<const double> b) {
    if (a.size() != b.size())
        throw std::invalid_argument("size mismatch");
    double sum = 0.0;
    for (std::size_t i = 0; i < a.size(); ++i)
        sum += a[i] * b[i];
    return sum;
}
\end{lstlisting}

Stosowanie \code{span} pozwala przyjmować
\code{vector}, \code{array} i tablice bez narzucania własności.
Jednocześnie nie wolno zachować widoku dłużej niż trwa jego bufor.

\subsection{Wartość czy referencja?}

Dla małych i tanio kopiowanych typów (np. \code{int}) przekazanie przez
wartość jest często najlepsze. Dla tekstu, który ma być zapisany
wewnątrz obiektu, idiom \code{std::string} przez wartość i
\code{std::move} do pola może uprościć przeciążenia konstruktorów.
Nie jest jednak bezwarunkowo najszybszy dla każdego rozkładu wejść:
dla samych \emph{lvalue} może wykonać dodatkowe kopiowanie.

\begin{lstlisting}[caption={Jeden konstruktor obslugujacy lvalue i rvalue (C++11).}]
#include <string>
#include <utility>

class Person {
    std::string name_;
public:
    explicit Person(std::string name)
        : name_(std::move(name)) {}
    const std::string& name() const noexcept { return name_; }
};
\end{lstlisting}

\section{Wyrażanie kontraktów w systemie typów}
\label{sec:cpp20-contracts}

Silny typ eliminuje klasę błędów polegających na pomyleniu wartości o
tej samej reprezentacji. \code{explicit} blokuje niepożądaną konwersję
niejawną. \code{[[nodiscard]]} pomaga wykryć ignorowanie wyniku;
nie gwarantuje jednak obowiązkowego ostrzeżenia we wszystkich trybach.

\begin{lstlisting}[caption={Strong types rozrozniajace dwa identyfikatory (C++20).}]
#include <cstdint>

struct UserId {
    std::uint64_t value;
    explicit constexpr UserId(std::uint64_t v) : value(v) {}
    friend constexpr bool operator==(UserId, UserId) = default;
};

struct OrderId {
    std::uint64_t value;
    explicit constexpr OrderId(std::uint64_t v) : value(v) {}
    friend constexpr bool operator==(OrderId, OrderId) = default;
};

[[nodiscard]] constexpr bool valid(UserId id) noexcept {
    return id.value != 0;
}

static_assert(valid(UserId{1}));
// valid(OrderId{1}); // blad kompilacji
\end{lstlisting}

\subsection{Concepts: wymóg składniowy nie zawsze jest wymogiem semantycznym}

\code{requires} pozwala sprawdzić, czy wyrażenie ma poprawną składnię,
ale nie dowodzi prawa matematycznego, np. przechodniości porządku,
asocjatywności czy stabilności haszowania. Dokumentacja musi dopowiadać
wymagania semantyczne. Przykładowo komparator dla sortowania musi
realizować \emph{strict weak ordering}; przypadkowy predykat typu
\code{bool(const T&,const T&)} nie wystarcza.

\begin{lstlisting}[caption={Ograniczenie generycznego sumowania (C++20).}]
#include <concepts>
#include <ranges>
#include <type_traits>

template<class R>
concept SummableRange = std::ranges::input_range<R> &&
    std::default_initializable<std::ranges::range_value_t<R>> &&
    requires(std::ranges::range_value_t<R>& result,
             std::ranges::range_reference_t<R> value) {
        result = result + value;
    };

template<SummableRange R>
auto sum_values(R&& range) {
    using T = std::ranges::range_value_t<R>;
    T result{};
    for (auto&& value : range) result = result + value;
    return result;
}
\end{lstlisting}

Wymóg \code{default_initializable} odpowiada instrukcji \code{T result\{\}}.
Wyrażenie w \code{requires} sprawdza dokładnie przypisanie wykonywane
w pętli, także gdy element zakresu jest typem pośrednim, a nie \code{T&}.
To nadal nie dowodzi sensu matematycznego sumowania. Jeśli wartość
domyślna nie jest elementem neutralnym dla dodawania, funkcja powinna
przyjmować wartość początkową od wywołującego.

\section{Obsługa błędów: wyjątki, typ wyniku i status}

\begin{center}\small
\begin{tabularx}{\textwidth}{@{}p{0.33\textwidth}X@{}}
\toprule
Kontrakt API & Kiedy i jak go użyć \\
\midrule
Wyjątek & Przerywa zwykły tok działania i przenosi błąd do obsługi; zasoby powinny być chronione przez RAII. \\
\code{std::optional<T>} & Oznacza obecność lub brak wyniku, gdy przyczyna braku nie jest potrzebna. \\
\code{std::expected<T,E>} & Zwraca wynik albo jawny błąd typu \code{E}; standard od C++23. \\
\code{std::error_code} & Przenosi kod błędu bez wyjątku, często obok wariantu zgłaszającego wyjątki. \\
\code{noexcept} & Jest częścią kontraktu: wyjątek opuszczający taką funkcję kończy program. \\
\bottomrule
\end{tabularx}
\end{center}

W C++20 podstawowymi wyborami są wyjątki oraz jawne typy wyniku,
np. \code{optional<T>} dla naturalnego braku wartości. Od C++23
\code{expected<T,E>} reprezentuje wartość lub opis błędu. Żadna
strategia nie jest uniwersalna: trzeba ustalić politykę dla całej
warstwy biblioteki, a nie mieszać bez powodu sposoby sygnalizowania.

\begin{lstlisting}[caption={Jawny blad jako typ wyniku w C++20.}]
#include <optional>
#include <string_view>

std::optional<unsigned> parse_digit(std::string_view x) {
    if (x.size() != 1 || x[0] < '0' || x[0] > '9')
        return std::nullopt;
    return static_cast<unsigned>(x[0] - '0');
}
\end{lstlisting}

\begin{lstlisting}[caption={Rozbudowany blad jako expected (C++23; fragment).}]
#include <cstddef>
#include <expected>
#include <string>
#include <string_view>

std::expected<std::size_t, std::string>
nonempty_length(std::string_view text) {
    if (text.empty())
        return std::unexpected<std::string>{"empty text"};
    return text.size();
}
\end{lstlisting}

\subsection{Gwarancje wyjątkowe i \texttt{noexcept}}

Funkcja modyfikująca stan może gwarantować: \emph{brak skutków po
niepowodzeniu} (silna gwarancja), zachowanie niezmienników mimo częściowej
zmiany (podstawowa gwarancja) albo brak wyjątków (\code{noexcept}).
\code{noexcept} nie znaczy, że błąd jest niemożliwy: nieobsłużony
wyjątek kończy wtedy program przez \code{std::terminate}.

\begin{lstlisting}[caption={Silna gwarancja przez przygotowanie kopii.}]
#include <string>
#include <utility>

class Settings {
    std::string data_;
public:
    void update(std::string next) {
        // Wszystkie kosztowne operacje wykonano przy
        // tworzeniu argumentu 'next'. Swap string ma noexcept.
        data_.swap(next);
    }
    const std::string& get() const noexcept { return data_; }
};
\end{lstlisting}

\section{Krycie implementacji: PImpl i granica ABI}
\label{sec:cpp20-pimpl}

Nagłówki udostępniające wszystkie pola klas narzucają zależności
kompilacyjne, a zmiana ich układu może naruszać ABI bibliotek
współdzielonych. Idiom PImpl (\emph{pointer to implementation})
przenosi szczegóły do osobnego typu. Cena: alokacja, pośredni dostęp
i potrzeba poprawnego zarządzania funkcjami specjalnymi.

\begin{lstlisting}[caption={widget.hpp -- publiczny naglowek.}]
#pragma once
#include <memory>
#include <string>

class Widget {
public:
    explicit Widget(std::string name);
    ~Widget();
    Widget(Widget&&) noexcept;
    Widget& operator=(Widget&&) noexcept;
    Widget(const Widget&) = delete;
    Widget& operator=(const Widget&) = delete;
    std::string name() const;
private:
    struct Impl;
    std::unique_ptr<Impl> impl_;
};
\end{lstlisting}

\begin{lstlisting}[caption={widget.cpp -- pelna definicja ukrytego typu.}]
#include "widget.hpp"
#include <utility>

struct Widget::Impl {
    std::string name;
    explicit Impl(std::string n) : name(std::move(n)) {}
};

Widget::Widget(std::string name)
    : impl_(std::make_unique<Impl>(std::move(name))) {}
Widget::~Widget() = default;
Widget::Widget(Widget&&) noexcept = default;
Widget& Widget::operator=(Widget&&) noexcept = default;
std::string Widget::name() const { return impl_->name; }
\end{lstlisting}

\begin{lstlisting}[caption={main.cpp -- uzycie niezalezne od szczegolow Impl.}]
#include "widget.hpp"
#include <iostream>

int main() {
    Widget w("sample");
    std::cout << w.name() << '\n';
}
\end{lstlisting}

Destruktor \code{Widget} jest definiowany \emph{po} pełnej definicji
\code{Impl}. Jest to istotne dla destrukcji \code{unique_ptr<Impl>}.
Wersja z domyślnym przenoszeniem dopuszcza obiekt \code{Widget}
w stanie \emph{moved-from} z pustym \code{impl_}; użytkownik nie może
wtedy wywoływać \code{name()} bez dodatkowo określonego kontraktu.
Jeśli wymagane jest, by każdy \code{Widget} zawsze miał implementację,
należy inaczej zaprojektować przenoszenie.

\subsection{Stabilność ABI nie jest automatyczna}

PImpl stabilizuje \emph{układ prywatnego stanu}, ale samo w sobie nie
zapewnia zgodności między kompilatorami, standardowymi bibliotekami,
wersjami runtime ani opcjami kompilacji. Publiczne typy
\code{std::string} lub \code{std::vector} mogą również wpływać na ABI.
Na granicy modułów współdzielonych wymagających zgodności wielu
środowisk warto rozważyć wąski interfejs C: nieprzezroczysty uchwyt,
stałe formaty danych i jawne funkcje tworzenia/niszczenia.

\section{Polimorfizm: dynamiczny, statyczny i wymazywanie typów}

Do rozszerzalnego API można zastosować wirtualną klasę bazową,
typ generyczny oparty na concepts, \code{variant} dla zamkniętego
zbioru alternatyw lub \emph{type erasure}. Mechanizmy różnią się
kosztami, możliwością rozbudowy i wymaganiami własności.

\begin{lstlisting}[caption={Interfejs dynamiczny i bezpieczne niszczenie polimorficzne.}]
#include <memory>
#include <string>

struct Renderer {
    virtual ~Renderer() = default;
    virtual std::string render(int value) const = 0;
};

struct SimpleRenderer final : Renderer {
    std::string render(int value) const override {
        return std::to_string(value);
    }
};

std::unique_ptr<Renderer> make_renderer() {
    return std::make_unique<SimpleRenderer>();
}
\end{lstlisting}

\begin{center}\small
\begin{tabularx}{\textwidth}{@{}p{0.25\textwidth}X@{}}
\toprule
Mechanizm & Atuty i koszty \\
\midrule
\code{virtual} & Otwarta hierarchia, wywołania pośrednie, ABI wrażliwe na zmianę interfejsu. \\
\code{template}/concept & Weryfikacja i specjalizacja przy kompilacji, więcej instancji kodu. \\
\code{variant} & Zamknięta lista typów, jawna obsługa przypadków, możliwy wzrost rozmiaru obiektu. \\
Type erasure & Jednolity typ opakowania, ukryta implementacja, zwykle pośredni dispatch i ewentualna alokacja. \\
\bottomrule
\end{tabularx}
\end{center}

\section{Punkty dostosowania: ADL, hidden friends i CPO}
\label{sec:cpp20-adl}

\emph{Argument-dependent lookup} (ADL) wyszukuje wolne funkcje
w przestrzeniach nazw związanych z typami argumentów. Pozwala
implementować operacje jak \code{swap}, \code{begin} czy
\code{operator==} bez modyfikowania biblioteki głównej.
\emph{Hidden friend} jest funkcją zaprzyjaźnioną zdefiniowaną
wewnątrz klasy; można ją odnaleźć przez ADL, ale nie przez zwykłe
wyszukiwanie nazwy bez odpowiedniego argumentu.

\begin{lstlisting}[caption={Hidden friend i idiom wywolywania swap przez ADL.}]
#include <utility>

namespace geom {
struct Point {
    int x{};
    int y{};
    friend bool operator==(const Point&, const Point&) = default;
    friend void swap(Point& a, Point& b) noexcept {
        using std::swap;
        swap(a.x, b.x);
        swap(a.y, b.y);
    }
};
}

template<class T>
void exchange(T& a, T& b) {
    using std::swap;
    swap(a, b); // moze znalezc geom::swap przez ADL
}
\end{lstlisting}

Nowoczesna biblioteka Ranges stosuje \emph{customization point objects}
(CPO), np. \code{std::ranges::begin}. Nie należy utożsamiać ich
z dowolnymi wolnymi funkcjami; obiekt CPO ma własne reguły wywołania,
weryfikacji wymagań i wyszukiwania implementacji.

\subsection{Wzorzec \texttt{tag\_invoke} -- nie standardowe API}

\code{tag_invoke} jest popularną konwencją proponowaną dla rozszerzalnych
CPO, ale \emph{nie jest standardową funkcją C++20 ani C++23}.
W bibliotekach, które ją przyjmują, użytkownik implementuje
\code{tag_invoke(tag, arguments...)} w przestrzeni nazw swego typu,
a biblioteka znajduje ją za pośrednictwem ADL. Bez własnej implementacji
mechanizmu nie wolno pisać \code{std::tag_invoke}.

\begin{lstlisting}[caption={Minimalna ilustracja konwencji tag\_invoke (fragment).}]
struct serialize_tag {};

namespace model {
struct Record { int id; };

// Odnajdywane przez ADL, gdy argumentem jest Record.
int tag_invoke(serialize_tag, const Record& r) {
    return r.id;
}
}

int serialize(const model::Record& r) {
    return tag_invoke(serialize_tag{}, r);
}
\end{lstlisting}

\section{Policy-based design i wstrzykiwanie zależności}

W \emph{policy-based design} strategia jest parametrem typu.
Wybrana implementacja może zostać wbudowana przez kompilator,
ale zwiększa liczbę instancji szablonów i zależności od nagłówków.
W \emph{dependency injection} obiekt otrzymuje zależności w konstruktorze
lub jako argument, zamiast samodzielnie tworzyć globalne usługi.

\begin{lstlisting}[caption={Polityka zapisu z kontrola w czasie kompilacji (C++20).}]
#include <concepts>
#include <string_view>

struct NullSink {
    void write(std::string_view) const noexcept {}
};

struct CounterSink {
    int* counter;
    void write(std::string_view) const noexcept { ++*counter; }
};

template<class T>
concept Sink = requires(T s, std::string_view text) {
    { s.write(text) } -> std::same_as<void>;
};

template<Sink S>
class Logger {
    S sink_;
public:
    explicit Logger(S sink) : sink_(sink) {}
    void log(std::string_view text) { sink_.write(text); }
};
\end{lstlisting}

Ograniczenie kompilacyjne nie zastępuje dokumentacji: każda polityka
\code{Sink} musi sama zapewniać odpowiednie własności współbieżne,
a \code{Logger} nie powinien obiecywać większej odporności na błędy
niż implementacja przekazanej polityki.
W pokazanym \code{CounterSink} wskaźnik musi pozostać ważny przez czas
używania loggera; dwa wątki nie mogą jednocześnie zwiększać tego samego
zwykłego licznika bez synchronizacji.

\section{Bezpieczeństwo wątkowe i reentrancy}

\emph{Thread-safe} nie jest jednoznacznym stwierdzeniem bez określenia
zakresu. Dokumentacja powinna mówić, czy dwa wątki mogą jednocześnie
wywoływać metody na tym samym obiekcie, czy tylko na różnych obiektach,
oraz czy zezwala się na jednoczesne czytanie i modyfikowanie.
\emph{Reentrant} oznacza, że wywołanie może bezpiecznie nakładać się
z kolejnym wywołaniem tej samej funkcji w określonym modelu; nie jest
automatycznym synonimem pełnej \emph{thread safety}.

Warto jawnie określić, czy callback dostarczony przez użytkownika
jest uruchamiany pod blokadą. Wywołanie obcego kodu przy zablokowanym
mutexie może spowodować deadlock, rekurencję lub trudny do przewidzenia
czas trzymania blokady.

\section{Publiczny nagłówek, moduły i konfiguracja CMake}

Czytelna biblioteka oddziela pliki \code{include/}, \code{src/},
\code{tests/} i \code{examples/}. W publicznym nagłówku umieszczamy
minimalny zbiór deklaracji potrzebnych odbiorcy, a nie \code{using
namespace std;} ani zależności od szczegółów implementacji.

\begin{lstlisting}[caption={Przykladowe CMakeLists.txt dla biblioteki publicznej.},language={}]
cmake_minimum_required(VERSION 3.20)
project(geometry VERSION 1.0 LANGUAGES CXX)

add_library(geometry src/geometry.cpp)
target_compile_features(geometry PUBLIC cxx_std_20)
target_include_directories(geometry PUBLIC
    $<BUILD_INTERFACE:${CMAKE_CURRENT_SOURCE_DIR}/include>
    $<INSTALL_INTERFACE:include>)

add_executable(example examples/main.cpp)
target_link_libraries(example PRIVATE geometry)
\end{lstlisting}

\nota{Deklaracje modułów C++20 mogą zmniejszyć zależność od klasycznego
preprocesora, ale same w sobie nie zapewniają stabilności ABI.
Konfiguracja modułów i dostępność \code{import std;} zależą od
wersji narzędzi; szczegóły omówiono w rozdziale~\ref{chap:cpp17-build-abi-modules}.}

\section{Wersjonowanie i testowanie kontraktu}

Zmiana kodu wewnętrznego nie jest zmianą API, dopóki zachowuje
zewnętrznie obiecane zachowanie. Zmiana przeciążenia, domyślnego
argumentu, wyjątków, czasu życia zwracanych referencji lub układu
publicznej klasy może mieć znacznie większe konsekwencje niż sugeruje
rozmiar poprawki. Warto odróżniać zgodność źródłową, ABI oraz zgodność
semantyczną.

\begin{enumerate}
\item Testuj kontrakt na granicy publicznych interfejsów, a nie wyłącznie
      metody prywatne; dodawaj testy przypadków błędnych.
\item Kompiluj testowego klienta w osobnym projekcie, korzystając tylko
      z publicznych nagłówków lub eksportowanych modułów.
\item Sprawdzaj zachowanie dla lvalue, rvalue, typów const oraz wyjątków
      we wszystkich operacjach z własnością zasobów.
\item Jeśli biblioteka obiecuje ABI, testuj zmianę wersji z programem
      zlinkowanym względem poprzedniego wydania.
\item Mierz koszt użycia biblioteki: narzut alokacji, wielkość binarium,
      opóźnienia i zachowanie w wielu wątkach.
\end{enumerate}

\section{Studium przypadku: API biblioteki numerycznej}

Załóżmy, że biblioteka oblicza iloczyn skalarny i przechowuje operator
liniowy. Dobry interfejs rozdziela widok danych od stanu własnego,
kontroluje zgodność wymiarów i określa zachowanie dla wyjątków.

\begin{lstlisting}[caption={Niewlasnosciowe wejscie do funkcji numerycznej (C++20).}]
#include <span>
#include <stdexcept>

[[nodiscard]] double dot(std::span<const double> lhs,
                         std::span<const double> rhs) {
    if (lhs.size() != rhs.size())
        throw std::invalid_argument("dimension mismatch");
    double acc = 0;
    for (std::size_t i = 0; i != lhs.size(); ++i)
        acc += lhs[i] * rhs[i];
    return acc;
}
\end{lstlisting}

Przy zastosowaniach wymagających powtarzalności numerycznej trzeba
udokumentować metodę sumowania i ewentualne przestawianie operacji.
Implementacja równoległa nie musi dać bitowo identycznego wyniku jak
sekwencyjna. Jeżeli biblioteka ma wspierać akceleratory i macierze
wielowymiarowe, należy oddzielić interfejs algorytmu od szczegółów
alokacji i układu pamięci.

\section{C API jako granica bibliotek współdzielonych}

Kiedy biblioteka ma być używana przez różne języki lub kompilatory,
wąski interfejs C bywa stabilniejszą granicą niż publiczna klasa C++.
Użytkownik nie musi znać rozmiaru obiektu, a biblioteka zachowuje
pełną kontrolę nad alokacją i zniszczeniem. Dwa podstawowe warunki to
\emph{nieprzekraczanie granicy przez wyjątki C++} oraz jawny kontrakt
zwalniania każdego zasobu.

\begin{lstlisting}[caption={calc\_api.h -- przyklad nieprzezroczystego uchwytu.}]
#pragma once
#ifdef __cplusplus
extern "C" {
#endif

typedef struct calc_handle calc_handle;
calc_handle* calc_create(int initial_value);
int calc_get(const calc_handle* handle, int* output);
void calc_destroy(calc_handle* handle);

#ifdef __cplusplus
}
#endif
\end{lstlisting}

\begin{lstlisting}[caption={calc\_api.cpp -- implementacja z obsluga bledu alokacji.}]
#include "calc_api.h"
#include <new>

struct calc_handle { int value; };

extern "C" calc_handle* calc_create(int value) {
    return new (std::nothrow) calc_handle{value};
}
extern "C" int calc_get(const calc_handle* h, int* out) {
    if (!h || !out) return 1;
    *out = h->value;
    return 0;
}
extern "C" void calc_destroy(calc_handle* h) {
    delete h;
}
\end{lstlisting}

Interfejs jest ilustracyjny. W realnym systemie warto wprowadzić
kody błędów, wersję ABI, zasady bezpieczeństwa wątkowego i politykę
alokacji. Nie należy zakładać, że wyjątki C++ mogą bezpiecznie
przechodzić przez granicę C lub modułu skompilowanego innym
kompilatorem. Szczególnie ważne jest, aby obiekt utworzony przez
\code{calc_create} zawsze niszczyć przez \code{calc_destroy},
a nie przez \code{free} ani \code{delete} po stronie klienta.

\section{Projektowanie wydajności: alokatory, PMR i batching}

Interfejs może narzucać niepotrzebne alokacje przez zwracanie
osobnych kontenerów dla każdej małej operacji. Często warto
rozważyć wariant zapisujący do bufora dostarczonego przez
wywołującego albo przetwarzanie wsadowe (\emph{batching}). Nie znaczy
to jednak, że każde API powinno przyjmować ręcznie zarządzane wskaźniki:
lepsza wydajność musi być uzasadniona pomiarem.

Biblioteka \code{std::pmr} pozwala oddzielić politykę alokacji od
typu danych. Przyjmowanie \code{std::pmr::memory_resource*} oznacza
jednak dodatkowe zobowiązanie: zasób pamięci musi istnieć co najmniej
tak długo, jak kontenery wykorzystujące go do zwalniania pamięci.

\begin{lstlisting}[caption={Kontener PMR z buforem przydzielonym przez klienta (C++17).}]
#include <array>
#include <cstddef>
#include <memory_resource>
#include <vector>

int main() {
    std::array<std::byte, 2048> storage{};
    std::pmr::monotonic_buffer_resource pool(
        storage.data(), storage.size());
    std::pmr::vector<int> values(&pool);
    for (int i = 0; i < 100; ++i) values.push_back(i);
}
\end{lstlisting}

\code{monotonic_buffer_resource} najpierw korzysta z tablicy
\code{storage}, ale po jej zapełnieniu może pobrać dodatkową pamięć
z zasobu nadrzędnego. Sama tablica nie narzuca więc twardego limitu.

\section{Dokumentacja API na przykładzie kontraktu operacji}

Dla każdej publicznej funkcji warto zapisać krótki, ale precyzyjny
kontrakt. W przypadku funkcji \code{dot(span,span)} może on brzmieć:

\begin{description}[style=nextline]
\item[Warunki wstępne] Obydwa widoki wskazują istniejące, czytelne
bufory, których czas życia obejmuje całe wywołanie.
\item[Wynik] Iloczyn skalarny w arytmetyce \code{double}.
\item[Błędy] Dla różnych długości zgłaszany jest
\code{std::invalid_argument}.
\item[Wyjątki] Przy błędzie argumentów żaden bufor nie jest
modyfikowany; sama funkcja nie przejmuje własności.
\item[Współbieżność] Równoczesne wywołania są dozwolone,
jeśli inni uczestnicy nie modyfikują tych samych elementów
w sposób powodujący wyścig danych.
\item[Wydajność] Koszt liniowy względem liczby elementów,
bez konieczności dodatkowej alokacji.
\end{description}

Umowa nie powinna obiecywać braku zależności od kolejności dodawania
dla liczb zmiennoprzecinkowych. Jeżeli kod ma gwarantować określony
błąd numeryczny, taki warunek trzeba jawnie sformułować i uzasadnić.

\section{Ewolucja API i polityka kompatybilności}

Podczas wydawania kolejnej wersji warto rozważyć odrębnie:
\emph{source compatibility}, \emph{binary compatibility} i
\emph{behavioral compatibility}. Na przykład dodanie publicznego
przeciążenia może prowadzić do niejednoznaczności wcześniej
poprawnego wywołania, mimo że żadna stara sygnatura nie została
usunięta. Zmiana wartości domyślnego argumentu może z kolei
prowadzić do innego zachowania po kompilacji klienta, choć ABI
pozostaje bez zmian.

W projektach opartych na bibliotekach współdzielonych stosuje się
niekiedy \emph{symbol versioning} lub oddzielne numery wersji ABI.
Nie ma jednak jednego mechanizmu przenośnego między wszystkimi
platformami. Dla bibliotek publicznych szczególnie przydatne są
przykłady minimalnych klientów kompilowane w macierzy wersji oraz
testy odtwarzające zachowanie poprzedniej wersji.

\section{Lista kontrolna przeglądu biblioteki}

\begin{enumerate}
\item Czy publiczne nagłówki kompilują się samodzielnie bez przypadkowych
zależności pośrednich?
\item Czy własność, możliwość \code{nullptr}, czas życia oraz
odpowiedzialność za zwalnianie zasobów są jednoznaczne?
\item Czy koncepty odpowiadają wszystkim wymaganym operacjom kodu?
\item Czy zachowanie przy błędzie jest testowane, łącznie z wyjątkami
z konstruktorów i callbacków użytkownika?
\item Czy funkcje \code{noexcept} naprawdę nie przepuszczają wyjątków?
\item Czy zbędne alokacje i niejawne kopiowanie są zmierzone, a nie
tylko zakładane?
\item Czy API można przetestować z implementacją testową zależności,
bez uruchamiania sieci ani bazy danych?
\item Czy publiczne typy nie ujawniają szczegółów, które trzeba
będzie zmienić w przyszłości?
\item Czy dokumentacja dokładnie określa dopuszczalne wywołania
współbieżne?
\end{enumerate}


\section*{Podsumowanie rozdziału}
\addcontentsline{toc}{section}{Podsumowanie rozdziału}

Projektowanie bibliotek C++ polega przede wszystkim na precyzyjnym
wyrażeniu zobowiązań. Typy takie jak \code{span}, \code{unique_ptr},
\code{optional}, concepts i interfejsy wirtualne są narzędziami do
modelowania tych zobowiązań. PImpl pomaga ukryć implementację; ADL
i CPO pomagają projektować rozszerzalność; testy, wersjonowanie
oraz dokumentacja wyznaczają granice rzeczywistej niezawodności.
Najważniejsza zasada brzmi: nie obiecuj silniejszych gwarancji własności,
współbieżności, bezpieczeństwa wyjątkowego ani ABI, niż możesz
zapewnić i przetestować.

\section*{Literatura i dokumentacja}
\addcontentsline{toc}{section}{Literatura i dokumentacja}
\begingroup\sloppy
\begin{itemize}
\item ISO C++ Core Guidelines, \url{https://isocpp.github.io/CppCoreGuidelines/CppCoreGuidelines}.
\item \url{https://en.cppreference.com/w/cpp/language/adl} --- ADL i hidden friends.
\item \url{https://en.cppreference.com/w/cpp/ranges} --- customization point objects Ranges.
\item \url{https://en.cppreference.com/w/cpp/language/pimpl} --- idiom PImpl.
\item \url{https://cmake.org/cmake/help/latest/guide/tutorial/} --- CMake.
\item Klaus Iglberger, \emph{C++ Software Design}, O'Reilly, 2022.
\end{itemize}
\endgroup
